\clearpage
\section{Exact numerical certificates}\label{sec:numeric}

The program \texttt{code/verify\_numeric.py} computes the Bell recurrence
\eqref{eq:bell-binomial} through index $1000$, using integer arithmetic only.
Its output is \texttt{code/numeric\_validation.json}. The first values are:
\begin{center}
\begin{tabular}{r r@{\qquad}r r}
\hline
$n$&$B_n$&$n$&$B_n$\\
\hline
0&1&8&4140\\
1&1&9&21147\\
2&2&10&115975\\
3&5&11&678570\\
4&15&12&4213597\\
5&52&13&27644437\\
6&203&14&190899322\\
7&877&15&1382958545\\
\hline
\end{tabular}
\end{center}

For the decimal counts reported in the source, the checker verifies the
following exact integer inequalities; no logarithm or floating-point
approximation enters the checks.
\begin{center}
\begin{tabular}{r r l}
\hline
$n$&Decimal digits of $B_n$&Verified interval\\
\hline
1&1&$10^0\leq B_1<10^1$\\
10&6&$10^5\leq B_{10}<10^6$\\
100&116&$10^{115}\leq B_{100}<10^{116}$\\
1000&1928&$10^{1927}\leq B_{1000}<10^{1928}$\\
\hline
\end{tabular}
\end{center}
These intervals certify the digit counts. For compact reproducibility, the
JSON retains each value's first and last at most $32$ decimal digits and a
SHA-256 fingerprint of its full decimal expansion. The fingerprints are
identifiers; the recurrence and the integer inequalities are the mathematical
checks. No value $B_{10000}$ is computed or asserted here.

Among $0\leq n\leq13$, precisely $n=2,3,7,13$ give primes:
\[
 B_2=2,\qquad B_3=5,\qquad B_7=877,\qquad B_{13}=27644437.
\]
For these four values $v$, the integer square roots are respectively
$1,2,29,5257$. The checker verifies $r^2\leq v<(r+1)^2$ and that no integer
$2\leq d\leq r$ divides $v$. This proves primality: if $v=ab$ with
$a,b>1$, at least one factor is at most $r$. For every other index through
$13$, the output records a divisor or the fact that the value is less than
two. The checks make no primality assertion about later Bell numbers and no
claim about infinitely many Bell primes.

All these are fixed finite computation certificates. Their recurrence loops,
exact divisions, decimal comparisons, and trial-divisor searches have explicit
finite bounds and can be expanded into ground arithmetic derivations. The
Python program is not a generic proof checker, and successful execution is
not presented as a completed syntactic formalization in \ipra.
